\documentclass[12pt,a4paper]{article}
\usepackage{iftex}
\ifPDFTeX
  \usepackage[T2A]{fontenc}
  \usepackage[utf8]{inputenc}
\else
  \usepackage{fontspec}
  \setmainfont{DejaVu Serif}
\fi
\usepackage[russian]{babel}
\usepackage{geometry}
\geometry{a4paper,left=20mm,right=20mm,top=20mm,bottom=23mm}
\usepackage{amsmath,amssymb,mathtools}
\usepackage{array,booktabs,longtable}
\usepackage{xcolor}
\usepackage{tikz}
\usepackage[unicode,colorlinks=true,linkcolor=blue!55!black]{hyperref}
\setlength{\parindent}{0pt}
\setlength{\parskip}{7pt}
\setlength{\emergencystretch}{3em}
\renewcommand{\arraystretch}{1.25}
\newcommand{\reviewlabel}[1]{\medskip\noindent\textbf{#1}\par}
\begin{document}

% Титульный лист
\thispagestyle{empty}
\begin{center}
\vspace*{0.28\textheight}
{\huge\bfseries Ревью домашней работы\par}
\vspace{22mm}
{\large Автор: Лёвин Артём Александрович\par}
\vspace{5mm}
{\large 16 октября 2026 года\par}
\vfill
\begin{tikzpicture}
  \draw[blue!45!black,line width=0.7pt]
  (-3.1,0) .. controls (-1.5,0.55) and (-0.8,-0.55) .. (0,0)
            .. controls (0.8,0.55) and (1.5,-0.55) .. (3.1,0);
\end{tikzpicture}
\vspace{12mm}
\end{center}
\newpage

\section*{Сводная таблица}
\small
\setlength{\tabcolsep}{3pt}
\begin{longtable}{@{}p{1.25cm}p{2.45cm}p{4.15cm}p{4.05cm}p{4.0cm}@{}}
\toprule
\textbf{№} & \textbf{Тема} & \textbf{Суть ошибки} & \textbf{Ваш ответ} & \textbf{Правильный ответ}\\
\midrule
\endhead
\hyperref[task:1]{1} & Включение границы & Левая граница ошибочно исключена. & $x\in(-1;+\infty)$ & $x\in[-1;+\infty)$\\
\addlinespace[3pt]
\hyperref[task:3]{3} & Запись луча & Правильная идея показана на оси, но требуемое неравенство не записано. & Рисунок; «не подходит» для $-3$ & $x>-3$; число $-3$ не принадлежит\\
\addlinespace[3pt]
\hyperref[task:6]{6} & Проверка чисел & Десятичная дробь $11{,}9$ не выбрана; вместо неё названо $9$. & «Только $9$ и $4$» & $4$ и $11{,}9$\\
\addlinespace[3pt]
\hyperref[task:8]{8} & Запись по рисунку & Промежуток записан, но в неравенстве пропущена переменная и не показана включённость левого конца. & $-3<1$; $x\in[-3;1)$ & $-3\le x<1$; $x\in[-3;1)$\\
\bottomrule
\end{longtable}
\normalsize
\newpage

\thispagestyle{empty}
\vspace*{\fill}
\begin{center}
{\LARGE\bfseries Разбор ошибок}
\end{center}
\vspace*{\fill}
\newpage

\phantomsection\label{task:1}
\section*{Задание № 1. Включение граничного значения}
\reviewlabel{Текст задания}
Запишите промежуток $x\ge -1$ с помощью знака принадлежности $\in$ и изобразите его на числовой оси.

\reviewlabel{Ваше решение}
На числовой оси отмечено число $-1$ пустым кружком, а направление продолжения промежутка показано вправо. Промежуток записан так:
\[
 x\in(-1;+\infty).
\]

\reviewlabel{Ваш ответ}
$\displaystyle x\in(-1;+\infty)$.

\reviewlabel{Ошибка}
\textcolor{red}{Ошибка:} левая граница неравенства исключена из промежутка, хотя по условию должна в него входить.

\reviewlabel{Где именно возникает ошибка}
Сначала Вы верно определили, что подходят числа справа от $-1$. Ошибка появляется при оформлении самой точки $-1$: пустой кружок и круглая скобка означают, что это число не входит в промежуток. Но условие содержит знак «больше или равно», поэтому точку исключать нельзя.

\reviewlabel{Почему это неверно}
Знак $\ge$ означает одновременное выполнение одного из двух случаев: число либо больше границы, либо равно ей. Следовательно, $x=-1$ допустимо. Для проверки достаточно подставить его в исходное неравенство:
\[
 -1\ge -1.
\]
Это верное утверждение: число равно самому себе. Запись же $x\in(-1;+\infty)$ эквивалентна более строгому неравенству $x>-1$ и не содержит $-1$. Таким образом, Ваш ответ описывает не тот набор чисел, который задан в условии.

\reviewlabel{Ключевая идея}
\textcolor{blue!60!black}{Ключевая идея:} знак строгого неравенства ($<$ или $>$) соответствует круглой скобке и пустой точке на оси; знак нестрогого неравенства ($\le$ или $\ge$) соответствует квадратной скобке и закрашенной точке. Бесконечность не является числом и всегда отделяется круглой скобкой.

\reviewlabel{Исправление от последнего правильного шага}
Сохраняем верно выбранное направление вправо. Заменяем пустую точку в $-1$ на закрашенную, а круглую левую скобку на квадратную:
\[
 x\in\boxed{[-1;+\infty)}.
\]

\reviewlabel{Полное правильное решение}
\begin{enumerate}
\item Из $x\ge-1$ следует, что подходят все числа больше $-1$, а также само число $-1$.
\item Значит, на числовой оси в точке $-1$ ставим закрашенный кружок.
\item От этой точки выделяем луч вправо, поскольку справа находятся числа, большие $-1$.
\item Левая квадратная скобка указывает, что граница включена. Правая круглая скобка у бесконечности указывает, что конечной правой границы нет.
\end{enumerate}
\[
\boxed{x\in[-1;+\infty).}
\]

\reviewlabel{Проверка}
Проверим число на границе и по одному числу с обеих сторон. Число $-1$ входит: $-1\ge-1$. Число $0$ входит: $0\ge-1$. Число $-2$ не входит: $-2\not\ge-1$. Полученный промежуток даёт именно такой результат.

\reviewlabel{Итог}
\textcolor{teal!60!black}{Итог:} требуется заменить левую круглую скобку на квадратную и сделать точку $-1$ на оси закрашенной: $x\in[-1;+\infty)$.

\reviewlabel{Задача на закрепление}
Запишите промежуток $x\ge-4$ с помощью знака $\in$ и изобразите его на числовой оси.
\newpage

\phantomsection\label{task:3}
\section*{Задание № 3. Неравенство по промежутку}
\reviewlabel{Текст задания}
Запишите неравенство, задающее промежуток $x\in(-3;+\infty)$. Укажите, подходит ли число $-3$.

\reviewlabel{Ваше решение}
Изображён луч, направленный вправо от $-3$, с незакрашенной точкой в $-3$. Для числа $-3$ записано: «не подходит». При этом требуемая буквенная запись неравенства отсутствует.

\reviewlabel{Ваш ответ}
Число $-3$ «не подходит»; неравенство не записано.

\reviewlabel{Ошибка}
\textcolor{red}{Ошибка:} ответ неполный: графическое изображение и верная проверка границы не заменяют неравенство, которое прямо требуется в задании.

\reviewlabel{Где именно возникает ошибка}
До вывода о границе ход рассуждения верен: Вы правильно показали, что $-3$ исключено, а справа от него лежат подходящие числа. Пропущен следующий обязательный шаг: перевести интервал с круглой левой скобкой в строгое неравенство $x>-3$.

\reviewlabel{Почему это неверно}
Запись $x\in(-3;+\infty)$ сообщает сразу две вещи. Во-первых, нижняя граница равна $-3$. Во-вторых, сама граница исключена: об этом говорит круглая скобка. Все числа из промежутка расположены правее $-3$, поэтому они \emph{строго больше} $-3$. Именно это свойство нужно выразить неравенством. Ответ «не подходит» раскрывает только судьбу одного числа, но не описывает все допустимые значения $x$.

\reviewlabel{Ключевая идея}
\textcolor{blue!60!black}{Ключевая идея:} запись промежутка и неравенство должны задавать одно и то же множество чисел. Для луча вправо от исключённой точки $a$ используется $x>a$; если точка включена, используется $x\ge a$.

\reviewlabel{Исправление от последнего правильного шага}
Сохраняем верную отметку открытой точки $-3$ и направление вправо. Добавляем недостающую запись:
\[
\boxed{x>-3.}
\]
Утверждение о том, что $-3$ не подходит, остаётся правильным.

\reviewlabel{Полное правильное решение}
\begin{enumerate}
\item Левая граница промежутка $(-3;+\infty)$ равна $-3$.
\item Круглая скобка около $-3$ означает, что равенство $x=-3$ не допускается.
\item Правая часть промежутка не ограничена; подходят все числа, которые больше $-3$.
\item Поэтому записываем строгое неравенство $x>-3$.
\item Отдельно проверяем число $-3$: строгое неравенство $-3>-3$ ложно, значит, число $-3$ промежутку не принадлежит.
\end{enumerate}
\[
\boxed{x>-3,\qquad -3\notin(-3;+\infty).}
\]

\reviewlabel{Проверка}
Число $-2$ находится правее $-3$ и удовлетворяет неравенству: $-2>-3$. Число $-4$ левее границы и не удовлетворяет неравенству. Число $-3$ также исключено, так как строгое равенство не допускается. Это совпадает с заданным промежутком.

\reviewlabel{Итог}
\textcolor{teal!60!black}{Итог:} построение выполнено по смыслу правильно, но к нему необходимо дописать строгое неравенство $x>-3$.

\reviewlabel{Задача на закрепление}
Запишите неравенство, задающее промежуток $x\in(2;+\infty)$. Укажите, принадлежит ли ему число $2$.
\newpage

\phantomsection\label{task:6}
\section*{Задание № 6. Числа внутри промежутка}
\reviewlabel{Текст задания}
Запишите в виде промежутка двойное неравенство $4\le x<12$. Какие из чисел $4$, $11{,}9$ и $12$ ему принадлежат?

\reviewlabel{Ваше решение}
Вы верно записали заданный промежуток:
\[
 x\in[4;12).
\]
Но при выборе чисел указали: «только $9$ и $4$».

\reviewlabel{Ваш ответ}
«Только $9$ и $4$».

\reviewlabel{Ошибка}
\textcolor{red}{Ошибка:} список подходящих чисел составлен неверно. В условии указана десятичная дробь $11{,}9$, а число $9$ отдельно в списке не дано. Кроме того, $11{,}9$ удовлетворяет двойному неравенству и должна быть выбрана.

\reviewlabel{Где именно возникает ошибка}
Последний полностью правильный шаг --- перевод $4\le x<12$ в $x\in[4;12)$. Первый ошибочный шаг --- вывод «только $9$ и $4$». При чтении списка важно различать запятую между числами и запятую как десятичный разделитель. Последовательность $11{,}9$ обозначает \emph{одно} число: одиннадцать целых и девять десятых.

\reviewlabel{Почему это неверно}
Принадлежность промежутку $[4;12)$ проверяется \emph{двумя} условиями одновременно:
\[
4\le a\qquad\text{и}\qquad a<12.
\]
Нельзя включать числа в ответ только потому, что они вообще находятся между $4$ и $12$: необходимо выбирать именно из предложенного списка. Даже если число $9$ удовлетворяет неравенству, это не делает его правильным вариантом ответа здесь, поскольку отдельно такого кандидата в условии нет. Проверка дроби $11{,}9$ показывает, что она больше $4$ и меньше $12$.

\reviewlabel{Ключевая идея}
\textcolor{blue!60!black}{Ключевая идея:} каждое предложенное значение нужно подставить в обе части двойного неравенства. Запись с запятой внутри числа, например $11{,}9$, означает десятичную дробь. Проверка принадлежности не зависит от того, целое число или дробное.

\reviewlabel{Исправление от последнего правильного шага}
Сохраняем верный промежуток $x\in[4;12)$. Последовательно проверяем три указанных в условии числа:
\begin{align*}
 4&:\quad 4\le 4<12 &&\text{--- верно};\\
 11{,}9&:\quad 4\le11{,}9<12 &&\text{--- верно};\\
 12&:\quad 4\le12<12 &&\text{--- неверно}.
\end{align*}
Таким образом, в ответе должны остаться $4$ и $11{,}9$.

\reviewlabel{Полное правильное решение}
\begin{enumerate}
\item Условие $4\le x$ допускает число $4$ и все большие него числа. Поэтому ставим слева квадратную скобку.
\item Условие $x<12$ не допускает число $12$. Поэтому справа ставим круглую скобку.
\item Получаем промежуток $x\in[4;12)$.
\item Число $4$ равно нижней границе, которая включена, поэтому подходит.
\item Число $11{,}9$ больше нижней границы и меньше верхней, поэтому подходит.
\item Число $12$ равно верхней границе, которая исключена, поэтому не подходит.
\end{enumerate}
\[
\boxed{x\in[4;12),\qquad\text{подходят }4\text{ и }11{,}9.}
\]

\reviewlabel{Проверка}
Число $11{,}9$ можно сравнить с верхней границей по десятым: $12=12{,}0$, а $11{,}9<12{,}0$. При этом $11{,}9>4$. Следовательно, дробь действительно принадлежит промежутку, а $12$ --- нет.

\reviewlabel{Итог}
\textcolor{teal!60!black}{Итог:} запись промежутка верна; ошибка только в выборе чисел. Правильный список: $4$ и $11{,}9$.

\reviewlabel{Задача на закрепление}
Запишите двойное неравенство $2\le x<9$ в виде промежутка. Какие из чисел $2$, $8{,}7$ и $9$ ему принадлежат?
\newpage

\phantomsection\label{task:8}
\section*{Задание № 8. Двойное неравенство по рисунку}
\reviewlabel{Текст задания}
По рисунку запишите двойное неравенство и промежуток. На числовой оси выделена часть от $-3$ до $1$: точка $-3$ закрашена, точка $1$ не закрашена.

\reviewlabel{Ваше решение}
Выписано неравенство
\[
 -3<1,
\]
а промежуток записан верно:
\[
 x\in[-3;1).
\]

\reviewlabel{Ваш ответ}
$\displaystyle -3<1$ и $\displaystyle x\in[-3;1)$.

\reviewlabel{Ошибка}
\textcolor{red}{Ошибка:} выражение $-3<1$ не является двойным неравенством для переменной $x$ и не задаёт выделенный промежуток. В нём отсутствует сама переменная и не отражено, что число $-3$ входит в промежуток.

\reviewlabel{Где именно возникает ошибка}
Вы правильно прочитали границы при записи $x\in[-3;1)$. Первый неверный шаг возникает при попытке записать двойное неравенство: вместо связи обеих границ с $x$ указано лишь истинное сравнение двух чисел, $-3<1$.

\reviewlabel{Почему это неверно}
Высказывание $-3<1$ верно независимо от того, какое значение имеет $x$. Поэтому оно не позволяет определить, подходит ли, например, число $0$, $-4$ или $2$. Двойное неравенство должно одновременно сообщать, что переменная \emph{не меньше} нижней границы и \emph{строго меньше} верхней. Здесь закрашенная точка $-3$ означает допустимость равенства слева, а пустая точка $1$ --- недопустимость равенства справа.

\reviewlabel{Ключевая идея}
\textcolor{blue!60!black}{Ключевая идея:} для ограниченного промежутка переменная располагается между двумя числами. Левая квадратная скобка превращается в знак $\le$, а правая круглая --- в знак $<$. Нужно записать обе связи с одной переменной: $a\le x<b$.

\reviewlabel{Исправление от последнего правильного шага}
Сохраняем правильно установленный промежуток $x\in[-3;1)$. Переносим его границы в двойное неравенство: нижняя включена, верхняя исключена.
\[
\boxed{-3\le x<1.}
\]

\reviewlabel{Полное правильное решение}
\begin{enumerate}
\item По рисунку выделены значения $x$ между числами $-3$ и $1$.
\item Левая точка $-3$ закрашена: разрешены и $-3$, и любые большие числа до правой границы. Значит, слева нужен знак $\le$.
\item Правая точка $1$ пустая: значение $1$ исключено. Значит, справа нужен знак $<$.
\item Ставим переменную между границами и получаем $-3\le x<1$.
\item Переводим запись в промежуток: квадратная скобка у $-3$ и круглая у $1$ дают $x\in[-3;1)$.
\end{enumerate}
\[
\boxed{-3\le x<1,\qquad x\in[-3;1).}
\]

\reviewlabel{Проверка}
Для $x=-3$ получаем $-3\le-3<1$: обе части верны, поэтому левая граница включена. Для $x=0$ получаем $-3\le0<1$: число находится внутри. Для $x=1$ получаем $-3\le1<1$: правая часть ложна, поэтому верхняя граница исключена. Для $x=-4$ нижнее условие тоже ложно. Все четыре проверки согласуются с рисунком.

\reviewlabel{Итог}
\textcolor{teal!60!black}{Итог:} запись промежутка правильная. В двойное неравенство нужно добавить переменную и верный знак у нижней границы: $-3\le x<1$.

\reviewlabel{Задача на закрепление}
На числовой оси выделены все числа от $-2$ включительно до $5$ не включительно. Запишите соответствующие двойное неравенство и промежуток.

\vfill
\begin{center}
\footnotesize Лёвин Артём Александрович эксклюзивно для Кирилла Зиновьева
\end{center}

\end{document}
